\documentclass[11pt,a4paper]{article}
\usepackage[T1]{fontenc}
\usepackage[utf8]{inputenc}
\usepackage[brazil]{babel}
\usepackage[margin=1.8cm]{geometry}
\usepackage[compact]{titlesec}
\titlespacing*{\section}{0pt}{0.9em}{0.4em}
\usepackage{amsmath,amssymb}
\usepackage{booktabs}
\usepackage{enumitem}
\usepackage{caption}
\usepackage{microtype}
\usepackage[hidelinks]{hyperref}

\captionsetup{font=small,labelfont=bf}
\setlist{nosep,leftmargin=1.4em}
\setlength{\parskip}{0.2em}
\pagestyle{plain}

\newcommand{\cmax}{C_{\max}}

\begin{document}

\noindent{\LARGE\bfseries Otimização do Sequenciamento na Produção de Baterias:\\[2pt] Abordagem via RCPSP}\\[6pt]
{\bfseries Etapa 5: Resultados do Método Exato}\\
Cristian Madeira (Pós), Otávio Novais (162962), Thaís Melo (164934)

\section*{Configuração dos experimentos}

O modelo de programação linear inteira mista (PLIM) tempo-indexado descrito na Etapa~4
\cite{pritsker} foi implementado em AMPL e resolvido com o Gurobi~13.0.3, acionado pela
biblioteca \texttt{amplpy}, com limite de 600~s por instância. Os experimentos foram executados
em um notebook com processador Intel Core i5 de 11ª geração e 16~GB de RAM, sob Linux.
Em relação à proposta da Etapa~4, foram adotados três refinamentos que não alteram o conjunto
de soluções ótimas:

\begin{itemize}
  \item \textbf{Precedência desagregada:} para cada arco $(i,j)$ e período $t$,
        $\sum_{s \le t} x_{js} \le \sum_{s \le t-p_i} x_{is}$ \cite{christofides}. Em testes
        preliminares (HiGHS, 60~s), deu limitantes inferiores melhores que a forma agregada em 8 das
        10 instâncias (iguais nas outras duas).
  \item \textbf{Horizonte e janelas:} $H$ é o menor makespan obtido pelo SSGS com sete regras de
        prioridade (LFT, LST, MSLK, MTS, GRPW, LPT, MIS) seguido de \emph{forward-backward
        improvement} \cite{valls}. As variáveis $x_{jt}$ só existem para $t \in [ES_j, LS_j]$,
        calculados pelo caminho crítico com esse $H$. Com o horizonte do arquivo ($H = 158$ em
        j301\_1), o modelo teria cerca de 5.000 binárias; com as janelas, tem entre 482 e 2.290.
  \item \textbf{Solução inicial:} o cronograma heurístico que define $H$ é fornecido ao solver
        como solução inicial.
\end{itemize}

Foram usadas dez instâncias do J30 da PSPLIB \cite{psplib}: j301\_1 (a instância apresentada na
Etapa~3, cujo arquivo-base no gerador se chama j30\_17) e nove instâncias dos grupos mais difíceis,
uma por grupo, todas com recursos escassos ($RS = 0{,}2$ ou $0{,}5$). Nos grupos com $RS \ge 0{,}7$,
o ótimo praticamente coincide com o limitante trivial (caminho crítico ou carga de recurso), e os
métodos não se diferenciam. Os ótimos de referência $C^*$ foram provados com o CP-SAT (OR-Tools). Todas as
soluções do Gurobi passaram por uma verificação independente do modelo (precedência, capacidade e makespan).

\section*{Resultados}

\begin{table}[ht]
\centering
\small
\setlength{\tabcolsep}{4.2pt}
\caption{Resultados do modelo PLIM tempo-indexado no Gurobi (limite de 600~s por instância).}
\label{tab:exato}
\begin{tabular}{lccrrrrrrrrcr}
\toprule
Instância & $RF$ & $RS$ & Var. & Restr. & $H$ & LP & $\cmax$ & LB & Gap (\%) & $C^*$ & Status & Tempo (s) \\
\midrule
% linhas geradas a partir de results/exact_selecao_j30_gurobi_D.csv
j301\_1 & 0,25 & 0,2 & 482 & 910 & 46 & 38,8 & \textbf{43} & 43 & 0,00 & 43 & ótimo & 0,08 \\
j309\_2 & 0,75 & 0,2 & 1.879 & 3.269 & 98 & 63,0 & \textbf{92} & 92 & 0,00 & 92 & ótimo & 286,5 \\
j3013\_1 & 1,00 & 0,2 & 1.061 & 1.865 & 64 & 48,3 & 59 & 53 & 10,17 & 58 & limite & 600,1 \\
j3025\_4 & 0,75 & 0,2 & 1.515 & 3.196 & 89 & 55,3 & \textbf{81} & 81 & 0,00 & 81 & ótimo & 242,8 \\
j3029\_8 & 1,00 & 0,2 & 1.454 & 3.043 & 89 & 66,5 & 81 & 72 & 11,11 & 80 & limite & 600,1 \\
j3037\_7 & 0,50 & 0,2 & 1.603 & 3.947 & 98 & 58,3 & \textbf{92} & 92 & 0,00 & 92 & ótimo & 40,9 \\
j3041\_10 & 0,75 & 0,2 & 1.740 & 4.253 & 106 & 69,4 & \textbf{99} & 99 & 0,00 & 99 & ótimo & 165,7 \\
j3045\_6 & 1,00 & 0,2 & 2.290 & 5.571 & 133 & 89,1 & \textbf{129} & 124 & 3,88 & 129 & limite & 600,1 \\
j3014\_2 & 1,00 & 0,5 & 599 & 1.126 & 56 & 46,8 & \textbf{53} & 53 & 0,00 & 53 & ótimo & 51,0 \\
j3030\_2 & 1,00 & 0,5 & 663 & 1.496 & 70 & 65,0 & \textbf{68} & 68 & 0,00 & 68 & ótimo & 6,0 \\
\midrule
Média & & & & & & & & & 2,52 & & 7/10 & 259,3 \\
\bottomrule
\end{tabular}
\par\smallskip
\begin{minipage}{0.97\textwidth}\footnotesize
$RF$, $RS$: fator e força de recurso (ProGen). Var./Restr.: variáveis binárias e restrições. $H$: horizonte (heurística). LP: relaxação linear.
$\cmax$: melhor solução inteira (em negrito quando igual a $C^*$). LB: melhor limitante inferior.
Gap $= (\cmax - \text{LB})/\cmax$. $C^*$: ótimo de referência. Status: ótimo provado ou limite de tempo.
\end{minipage}
\end{table}

\section*{Análise}

\textbf{Qualidade das soluções.} O Gurobi provou a otimalidade em 7 das 10 instâncias, com tempos
entre 0,08~s e 286~s (mediana de 51~s). Nas outras três, o limite de 600~s foi atingido, mas as
soluções continuam muito boas: $\cmax$ coincide com $C^*$ em 8 das 10 instâncias, e em j3013\_1 e
j3029\_8 fica uma unidade acima do ótimo (1,7\% e 1,3\%). Em j3045\_6 o solver encontrou o ótimo
(129), mas não conseguiu prová-lo: o limitante inferior parou em 124 (gap de 3,9\%). O desvio
médio em relação a $C^*$ é de 0,30\%, e o gap médio é de 2,5\% (8,4\% nas três instâncias não
resolvidas).

\textbf{Onde o modelo encontra dificuldade.} As três instâncias não resolvidas são exatamente as
dos grupos com $RF = 1{,}0$ e $RS = 0{,}2$: todas as atividades usam os quatro recursos, e esses
recursos são escassos. A relaxação linear fica, em média, 22\% abaixo do ótimo, mas esse
afastamento não explica sozinho a dificuldade: j3037\_7 (37\%) foi resolvida em 41~s, enquanto
j3013\_1 (17\%) não foi, o que indica que o fator de recurso pesa mais que a fraqueza da relaxação.

\textbf{Contribuição da heurística.} O horizonte $H$ do SSGS ficou, em média, 7\% acima do ótimo,
e o solver melhorou essa solução inicial em todas as instâncias (de 2 a 8 unidades de tempo).

\section*{Comparação preliminar com as metaheurísticas}

Para situar o método exato e responder à questão levantada em sala sobre o desempenho de
algoritmos genéticos neste problema, as mesmas instâncias foram resolvidas com três metaheurísticas
do framework RKO (\emph{Random-Key Optimizer}) \cite{rko}: o algoritmo genético por chaves
aleatórias (GA), o BRKGA \cite{brkga} e o RKO completo, que executa nove metaheurísticas em
paralelo. Os três usam o mesmo decodificador (chaves $\to$ lista de prioridade $\to$ SSGS, com ou
sem FBI), com 10 execuções por instância (10~s no J30; 30~s em cinco instâncias J60 com $RS = 0{,}2$),
na mesma máquina e com parada ao atingir o ótimo.

\begin{table}[ht]
\centering
\small
\caption{Metaheurísticas: RPD médio $= 100\,(\cmax - C^*)/C^*$ e execuções que atingiram $C^*$.}
\label{tab:mh}
\begin{tabular}{llrrr}
\toprule
Decodificador & Método & RPD J30 (\%) & RPD J60 (\%) & Execuções no ótimo \\
\midrule
SSGS + FBI & GA    & 0,29 & 0,92 & 94/150 \\
           & BRKGA & 0,24 & 1,00 & 95/150 \\
           & RKO   & 0,28 & 1,00 & 94/150 \\
\midrule
Só SSGS    & GA    & 0,43 & 2,00 & 79/150 \\
           & BRKGA & 0,51 & 2,07 & 70/150 \\
           & RKO   & \textbf{0,38} & \textbf{1,69} & \textbf{86/150} \\
\bottomrule
\end{tabular}
\end{table}

\vspace{-0.4em}
\textbf{Métodos de busca.} Com o FBI, os três métodos são estatisticamente equivalentes
(Friedman, $p = 0{,}81$): quando o GA usa a mesma representação e o mesmo decodificador, não se
observa a dificuldade esperada. Sem o FBI, o RKO obtém os melhores valores, mas a diferença não é
significativa após a correção para comparações múltiplas (Wilcoxon com Holm, $p = 0{,}13$). O fator
de maior efeito é o decodificador: no J60, o FBI reduz o RPD médio de 1,7--2,1\% para 0,9--1,0\% nos três
métodos ($p < 0{,}01$). Uma limitação é o orçamento: em 10--30~s cada execução avalia milhões de
cronogramas, muito acima dos 5.000 a 50.000 usados na literatura, o que tende a aproximar os métodos.

\textbf{Exato $\times$ metaheurística.} No J30, a melhor solução do RKO é igual ou melhor que a do
Gurobi nas dez instâncias, obtida em 0,58~s em média, contra 259~s do Gurobi. Em j3013\_1, o RKO
alcançou o ótimo (58) em 7 das 10 execuções, enquanto o Gurobi terminou em 59 após 600~s; em
j3029\_8, nenhum dos dois atingiu o ótimo (80). A vantagem do método exato está na prova de
otimalidade, obtida em 7 das 10 instâncias.

\textbf{Implicações.} O modelo exato é adequado para o J30, mas não garante a prova de otimalidade
em 600~s quando os recursos são mais disputados. O RKO será adotado como metaheurística do trabalho,
e a comparação será repetida com orçamento fixo de cronogramas, como na literatura.

\vspace{-0.6em}
{\scriptsize
\renewcommand{\refname}{\normalsize Referências}
\begin{thebibliography}{9}\setlength{\itemsep}{0pt}
\bibitem{christofides} Christofides, N. et al. Project scheduling with resource constraints: a branch and bound approach. \emph{European Journal of Operational Research}, 29(3):262--273, 1987.
\bibitem{psplib} Kolisch, R.; Sprecher, A. PSPLIB -- a project scheduling problem library. \emph{European Journal of Operational Research}, 96(1):205--216, 1997.
\bibitem{pritsker} Pritsker, A. A. B. et al. Multiproject scheduling with limited resources: a zero-one programming approach. \emph{Management Science}, 16(1):93--108, 1969.
\bibitem{rko} Chaves, A. A. et al. A random-key optimizer for combinatorial optimization. \emph{Journal of Heuristics}, 31(4), 2025.
\bibitem{brkga} Gonçalves, J. F.; Resende, M. G. C. Biased random-key genetic algorithms for combinatorial optimization. \emph{Journal of Heuristics}, 17(5):487--525, 2011.
\bibitem{valls} Valls, V.; Ballestín, F.; Quintanilla, S. Justification and RCPSP: a technique that pays. \emph{European Journal of Operational Research}, 165(2):375--386, 2005.
\end{thebibliography}}

\end{document}
